% =====================================================================
% ACE: proved unrestricted residue mean; open charge-preserving estimate
% Lifts No-Go only when charge-preserving ACE + bridge are both supplied
% =====================================================================

\chapter{A Priori Charge-Correcting Estimate: Status and Open Problem}
\label{ch:ACE}

\section{Role relative to the No-Go}

The a priori charge-correcting estimate (ACE) is the missing ingredient that
would lift the No-Go of Chapter~\ref{ch:nogo-attractor}. Its required form is
a strict upper bound
\begin{equation}
\label{eq:ACE-req}
\mathbb{E}[\chi]
\;\le\;
-\chi_{\min}
\;<\;0
\end{equation}
on the expected logarithmic expansion factor per step, obtained \textbf{solely}
from
\begin{enumerate}[label=(\roman*)]
  \item the residue-class structure of the map, and
  \item the topology of the \(243\) towers,
\end{enumerate}
without any reference to the numbers \(539\), \(61\), or \(G_4\).

This chapter:
\begin{itemize}
  \item proves the unrestricted residue average (crude ACE, non-circular);
  \item defines the charge-preserving correction \(2\cdot 3^k\);
  \item records why that correction reopens the expectation;
  \item states the ACE as an \textbf{open requirement}, not an established lemma;
  \item notes that a bridge \(\Psi\) is still needed after any ACE.
\end{itemize}

% ---------------------------------------------------------------------
\section{Unrestricted residue averages (proved)}
\label{sec:unrestricted}

\begin{definition}[Leading ratios of ordinary \(T_3\)]
\label{def:ratios}
For the ordinary (unrestricted) ternary Syracuse map, the magnitude ratios
associated to the three residue classes are taken at leading order as
\begin{equation}
\label{eq:ratios}
\rho_0=\frac13,
\qquad
\rho_1=\frac43,
\qquad
\rho_2=\frac23,
\end{equation}
corresponding to the branches \(n/3\), \((4n+2)/3\), and \((2n+1)/3\) in the
large-\(n\) continuum approximation (exact integer branches as in the Global
Lemma Core / HQH-539 specification).
\end{definition}

\begin{definition}[Logarithmic expansion factor (unrestricted)]
\label{def:chi-unres}
Under the uniform residue measure \(\mathbb{P}(r)=1/3\) for \(r\in\{0,1,2\}\),
\begin{equation}
\label{eq:chi-unres}
\chi_{\mathrm{unres}}(r)
\;:=\;
\ln\rho_r,
\qquad
\mathbb{E}_{\mathrm{unres}}[\chi]
\;:=\;
\frac13\sum_{r=0}^{2}\ln\rho_r.
\end{equation}
\end{definition}

\begin{theorem}[Unrestricted a priori log-mean]
\label{thm:unres-ACE}
\begin{equation}
\boxed{
\mathbb{E}_{\mathrm{unres}}[\chi]
\;=\;
\frac13\Bigl(
\ln\frac13+\ln\frac43+\ln\frac23
\Bigr)
\;=\;
\frac13\ln\frac{8}{27}
\;=\;
\ln\Bigl(\bigl(\tfrac{8}{27}\bigr)^{1/3}\Bigr)
\;\approx\;
-0.4054651081
\;<\;0.
}
\label{eq:unres}
\end{equation}
Equivalently,
\[
\exp\bigl(\mathbb{E}_{\mathrm{unres}}[\chi]\bigr)
\;=\;
\bigl(\tfrac{8}{27}\bigr)^{1/3}
\;=\;
\frac23.
\]
This number is strictly negative and independent of every global scale in the
model (\(539\), \(61\), \(G_4\), \(N_{\mathrm{flux}}\), seed size).
\end{theorem}

\begin{proof}
\[
\ln\frac13+\ln\frac43+\ln\frac23
=
\ln\Bigl(\frac13\cdot\frac43\cdot\frac23\Bigr)
=
\ln\frac{8}{27}.
\]
Divide by \(3\). Since \(8/27<1\), the log-mean is negative. No global integer
enters the calculation.
\end{proof}

\begin{corollary}[Crude a priori contraction]
\label{cor:crude}
Theorem~\ref{thm:unres-ACE} already supplies a crude, a priori contraction rate
in the logarithmic metric: the uniform residue average of leading ratios is
strictly contractive on average. This observation is non-circular and survives
the No-Go as a qualitative lemma.
\end{corollary}

\begin{corollary}[Weakness of the crude bound]
\label{cor:weak}
The unrestricted mean does \textbf{not}:
\begin{enumerate}[label=(\roman*)]
  \item force the specific integer \(539\);
  \item select a unique set of generational windings \(w_j=539+61j\);
  \item control the charge-preserving dynamics of Section~\ref{sec:charge}.
\end{enumerate}
It is too weak for Banach depth selection or for a unique generational
dictionary.
\end{corollary}

% ---------------------------------------------------------------------
\section{Charge-preserving map and correcting exponent}
\label{sec:charge}

\begin{definition}[Charge-preservation condition]
\label{def:charge-pres}
A trajectory is \emph{charge-preserving} if a declared ternary / flux charge
(e.g.\ residual class modulo \(3\) or modulo \(9\), as fixed by the model’s
charge ledger) is restored after each step, so that the orbit remains inside a
specified charge-preserving subspace \(\mathcal{C}\subset\mathbb{N}\).
\end{definition}

\begin{definition}[Minimal correcting exponent]
\label{def:k-min}
In the charge-preserving version of the map, the ordinary branch that would
send \(n\equiv 2\pmod{3}\) to \((2n+1)/3\) is replaced by a corrected step that
adds a term \(2\cdot 3^k\), where
\begin{equation}
\label{eq:k-min}
k=k(n)
\;:=\;
\min\bigl\{
\kappa\in\mathbb{Z}_{\ge 0}
:
\text{the corrected image restores the required residue modulo \(9\)}
\bigr\}.
\end{equation}
Thus \(k\) is a \emph{local} function of \(n\). Write the corrected step
schematically as
\begin{equation}
\label{eq:T3-corr}
T_3^{\mathrm{cp}}(n)
\;=\;
\frac{2n+1+2\cdot 3^{k(n)}}{3}
\quad\text{(when the \(\equiv 2\) branch is charge-corrected)},
\end{equation}
with the other residue branches unmodified or similarly constrained as
required by the ledger (the essential point is the presence of a
state-dependent \(3^{k(n)}\) correction).
\end{definition}

\begin{definition}[Logarithmic expansion factor (charge-preserving)]
\label{def:chi-cp}
Along a charge-preserving step with the published rule
\(T(n,k)=(n+1)/3+2\cdot 3^{k}\),
\begin{equation}
\label{eq:chi-cp}
\chi_{\mathrm{cp}}(n)
\;:=\;
\ln\Bigl(\frac{T(n,k(n))}{n}\Bigr)
\quad(n\ge 1,\ k(n)\ \text{exists}).
\end{equation}
\end{definition}

\begin{theorem}[Distribution of \(k(n)\) under \(Q=n\bmod 9\)]
\label{thm:k-dist}
For the published correction family \(T(n,k)=(n+1)/3+2\cdot 3^{k}\) and
preservation \(T\equiv n\pmod{9}\):
\begin{enumerate}[label=(\roman*)]
  \item \(2\cdot 3^{k}\bmod 9\in\{2,6,0\}\) according as \(k=0\), \(k=1\), or
        \(k\ge 2\). Hence \(k\ge 3\) never improves on \(k=2\), and any existing
        minimal exponent satisfies \(k(n)\in\{0,1,2\}\).
  \item Existence and value of \(k(n)\) depend only on \(n\bmod 27\). Exactly
        three of the nine classes with \(n\equiv 2\pmod{3}\) are feasible:
        \[
        n\equiv 17\pmod{27}\Rightarrow k=0,\quad
        n\equiv 23\pmod{27}\Rightarrow k=1,\quad
        n\equiv 14\pmod{27}\Rightarrow k=2.
        \]
        The other six classes admit \textbf{no} preserving \(k\).
  \item Uniformly over branch-2 integers:
        \(P(\mathrm{feasible})=1/3\), \(P(\mathrm{impossible})=2/3\), and
        conditional on feasibility, \(k\) is uniform on \(\{0,1,2\}\) with
        \(\mathbb{E}[k\mid\mathrm{feas}]=1\).
  \item On every feasible ray, \(T/n\to 1/3\) and
        \(\chi_{\mathrm{cp}}\to\ln(1/3)<0\); there is no runaway from large \(k\).
\end{enumerate}
(Full tables and Monte Carlo: \texttt{k\_n\_Distribution\_Analysis.md},
\texttt{scripts/analyze\_k\_distribution.py}.)
\end{theorem}

\begin{lemma}[Expectation is not the unrestricted residue mean]
\label{lem:not-residue}
Under charge preservation, \(\mathbb{E}[\chi_{\mathrm{cp}}]\) is \textbf{not}
equal to \(\frac13\ln(8/27)\). On the feasible set the leading ratio is
\(1/3\) (not the unrestricted branch-2 factor \(2/3\)). Moreover the map is
undefined on a density-\(2/3\) subset of branch-2 integers unless a secondary
rule is supplied. Trajectory averages therefore require the completed dynamics
and its invariant measure, not residue probabilities alone.
\end{lemma}

\begin{remark}[Correction to the informal ACE difficulty]
The difficulty is \textbf{not} that \(k(n)\) becomes arbitrarily large (it does
not, under the published family). The difficulty is: (i)~impossible classes of
density \(2/3\); (ii)~the invariant measure of the completed map; (iii)~the
bridge from \(\mathbb{E}[\chi]\) to the model depth \(\sigma=539\)
(not to the ACE symbol \(N_\star\), which is reserved for the short e-fold depth).
\end{remark}

% ---------------------------------------------------------------------
\section{Available ingredients (insufficient as yet)}
\label{sec:ingredients}

\begin{definition}[Residue Markov chain modulo \(9\)]
\label{def:markov9}
Let \(M_9\) be the Markov chain on \(\mathbb{Z}/9\mathbb{Z}\) induced by the
charge-preserving map \(T_3^{\mathrm{cp}}\). Transition probabilities
\(P_{ij}\) are, in principle, functions of the residue rules and the
charge-restoration condition only---independent of \(539\), \(61\), and \(G_4\).
\end{definition}

\begin{definition}[Democratic tower average]
\label{def:dem-avg}
The flux-democratic average over the \(243=3^5\) identical towers, with
canonical seed multiplicities \(20\) and \(21\) as in the published tower
tallies, is
\begin{equation}
\label{eq:dem}
\langle F\rangle_{243}
\;=\;
\frac{1}{243}\sum_{\tau=1}^{243} F(\tau),
\end{equation}
again without reference to a global orbit length.
\end{definition}

\begin{lemma}[Markov data and \(k\)]
\label{lem:markov-insufficient}
Under Theorem~\ref{thm:k-dist}, \(k\) is already bounded and is a function of
\(n\bmod 27\). The remaining Markov task is not to bound \(k\), but to obtain
the \emph{stationary distribution} of the completed charge-preserving map
(including a rule on the density-\(2/3\) impossible classes) so that
\(\mathbb{E}_\pi[\chi]\) is well-defined.
\end{lemma}

\begin{lemma}[Democracy does not supply the stationary average]
\label{lem:dem-insufficient}
The democratic tower average fixes the ensemble of initial seeds (including
multiplicities \(20\) and \(21\)). It does not by itself define the dynamics on
the impossible classes, nor compute \(\mathbb{E}_\pi[\chi]\) for a stationary
measure \(\pi\) of a completed map.
\end{lemma}

\begin{theorem}[ACE sharpened, not solved]
\label{thm:no-ACE-yet}
The \(k(n)\) analysis (Theorem~\ref{thm:k-dist}) \textbf{eliminates runaway
expansion} but does \textbf{not} supply:
\begin{enumerate}[label=(\alph*)]
  \item a completion of the dynamics on the density-\(2/3\) impossible set;
  \item a stationary expectation \(\mathbb{E}_\pi[\chi]\) for the resulting map;
  \item a non-circular bridge from that mean to the integer \(539\).
\end{enumerate}
Until (a)--(b) are exhibited from residue Markov structure and the
\(243\)-tower average alone (no \(539\), \(61\), \(G_4\)), the ACE remains an
\textbf{open requirement} and the No-Go of Chapter~\ref{ch:nogo-attractor}
\textbf{stands}. Item (c) becomes relevant only after a strictly negative
stationary mean is established.
\end{theorem}

\begin{proof}
Theorem~\ref{thm:k-dist} settles boundedness of \(k\) and the feasible/impossible
partition. Lemmas~\ref{lem:not-residue}--\ref{lem:dem-insufficient} show that
the stationary mean is not yet defined. No published derivation produces
\(\mathbb{E}_\pi[\chi]\le -\chi_{\min}<0\) free of global orbit length. Hence the
No-Go is not lifted.
\end{proof}

% ---------------------------------------------------------------------
\section{Sharpened open problem (completed map \(\to\) stationary mean \(\to\) bridge)}
\label{sec:open}

\begin{definition}[Impossible set]
\label{def:imp-set}
\[
\mathcal{I}
\;:=\;
\bigl\{
n\equiv 2\pmod{3}
:
n\bmod 27\in\{2,5,8,11,20,26\}
\bigr\}.
\]
Density of \(\mathcal{I}\) among branch-2 integers is \(2/3\).
\end{definition}

\begin{definition}[Candidate completion rules (not yet selected)]
\label{def:completion}
A \emph{completion} of the charge-preserving dynamics is any rule defining
\(T^\sharp(n)\) for \(n\in\mathcal{I}\), for example:
\begin{enumerate}[label=(\roman*)]
  \item \textbf{Projection:} map \(n\) to the nearest feasible class
        (e.g.\ reduce modulo a forced lattice) then apply feasible \(k\);
  \item \textbf{Rejection:} refuse the branch-2 step; fall back to an ordinary
        unrestricted \(T_3\) step (charge ledger updated by a declared defect);
  \item \textbf{Alternative charge-preserving branch:} a different additive
        family (not \(2\cdot 3^{k}\)) that restores \(Q\) on \(\mathcal{I}\).
\end{enumerate}
The ACE requires that \emph{some} completion be fixed from residue/tower data
alone and then used consistently---not that all three be developed.
\end{definition}

\begin{openproblem}[Completed map]
\label{op:complete}
Exhibit an explicit completion \(T^\sharp\) on \(\mathcal{I}\)
(Definition~\ref{def:completion}) such that:
\begin{enumerate}[label=(\roman*)]
  \item \(T^\sharp\) is defined for all \(n\in\mathbb{N}\) of interest;
  \item its restriction to feasible branch-2 classes agrees with
        \(T(n,k(n))\) of Theorem~\ref{thm:k-dist};
  \item the rule is stated using only residue-class and \(243\)-tower data
        (no \(539\), \(61\), \(G_4\)).
\end{enumerate}
\end{openproblem}

\begin{openproblem}[Stationary ACE]
\label{op:ACE}
Given a completion \(T^\sharp\) (Open Problem~\ref{op:complete}), let \(\pi\) be
a stationary probability measure for the induced residue Markov structure
(e.g.\ on \(\mathbb{Z}/27\mathbb{Z}\) or the minimal abelian state space),
averaged democratically over the \(243\) towers with seeds of types \(20\) and
\(21\). Prove
\begin{equation}
\label{eq:open-ACE}
\mathbb{E}_\pi[\chi]
\;:=\;
\int \ln\Bigl(\frac{T^\sharp(n)}{n}\Bigr)\,d\pi
\;\le\;
-\chi_{\min}
\;<\;0
\end{equation}
with \(\chi_{\min}>0\) independent of any pre-selected global orbit length.
The expectation must be obtained from the residue Markov structure and the
\(243\)-tower average alone.
\end{openproblem}

\begin{openproblem}[Bridge after ACE]
\label{op:bridge}
\textbf{Only after} Open Problem~\ref{op:ACE} yields a strictly negative
stationary mean does the following become relevant: supply a separate bridge
\begin{equation}
\label{eq:bridge-ACE}
N_\star
\;=\;
\Psi\bigl(\chi_{\min},\,243,\,\text{seed multiplicities }20,21\bigr)
\end{equation}
to a candidate integer orbit length, still without inserting \(539\) into
\(\Psi\). The bridge is not a substitute for the stationary ACE.
\end{openproblem}

\begin{theorem}[Order of operations]
\label{thm:order}
The logical order is strict:
\begin{equation}
\boxed{
\begin{aligned}
&\text{(1) complete dynamics on }\mathcal{I}
\\
&\text{(2) stationary mean }\mathbb{E}_\pi[\chi]\le -\chi_{\min}<0
\\
&\text{(3) bridge }\Psi\to N_\star
\\
&\text{(4) only then: democracy}\Rightarrow\text{Banach depth / unique }w_j.
\end{aligned}
}
\end{equation}
Skipping (1)--(2) and invoking (3)--(4) re-enters the circular
contraction-factor schema forbidden by the No-Go.
\end{theorem}

% ---------------------------------------------------------------------
\section{Status table}
\label{sec:status}

\begin{center}
\begin{tabular}{@{}p{0.40\textwidth}p{0.24\textwidth}p{0.28\textwidth}@{}}
\toprule
Statement & Status & Depends on \(539\)? \\
\midrule
\(\mathbb{E}_{\mathrm{unres}}[\chi]=\frac13\ln(8/27)<0\) &
\textbf{Proved} &
No \\
\(k\in\{0,1,2\}\) or impossible; no runaway &
\textbf{Proved} (Thm.~\ref{thm:k-dist}) &
No \\
Completion of dynamics on \(\mathcal{I}\) (density \(2/3\)) &
\textbf{Open} (Op.~\ref{op:complete}) &
Must not \\
Stationary ACE \(\mathbb{E}_\pi[\chi]\le -\chi_{\min}<0\) &
\textbf{Open} (Op.~\ref{op:ACE}) &
Must not \\
Bridge \(\Psi\to N_\star\) &
\textbf{Open; deferred} until ACE &
Must not \\
Flux democracy lifts No-Go &
\textbf{Blocked}; No-Go stands &
--- \\
Empirical Resonant Attractor protocol &
\textbf{Allowed} (No-Go ch.) &
Do not insert \(539\) into definition \\
\bottomrule
\end{tabular}
\end{center}

% ---------------------------------------------------------------------
\section{Consequences for the No-Go and the Banach rate}
\label{sec:conseq}

\begin{theorem}[Only qualitative contraction survives rigorously]
\label{thm:only-qual}
In the absence of the charge-preserving ACE
(Open Problem~\ref{op:ACE}), the only rigorous contraction statement that
survives as an a priori, non-circular lemma is
Corollary~\ref{cor:crude}: the uniform average over the three ordinary residue
classes is already negative,
\[
\frac13\ln\frac{8}{27}<0.
\]
That observation is too weak to force the specific integer \(539\) or to select
a unique set of generational windings.
\end{theorem}

\begin{corollary}[Banach rate \(\ln 3/539\) remains conditional]
\label{cor:banach-cond}
The Global Lemma Core identity \(\lambda=\ln 3/539<1\) remains a
\textbf{conditional} statement given \(\sigma=539\), not a consequence of
flux democracy or of Theorem~\ref{thm:unres-ACE} alone.
\end{corollary}

\begin{corollary}[No-Go unchanged]
\label{cor:nogo-stands}
Chapter~\ref{ch:nogo-attractor} stands: until Open Problems~\ref{op:ACE} and
\ref{op:bridge} are solved, flux democracy cannot be invoked to break
contraction circularity or to select a unique \(\sigma\)-dependent generational
dictionary.
\end{corollary}

% ---------------------------------------------------------------------
\section{Summary chain}
\label{sec:ace-summary}

\begin{equation}
\boxed{
\begin{aligned}
&\text{Unrestricted: }
\tfrac13\ln\tfrac{8}{27}<0
\quad\text{\textbf{proved}}
\\
&\text{\(k\)-law: bounded, period \(27\); no runaway}
\quad\text{\textbf{proved}}
\\
&\text{Complete map on density-\(2/3\) }\mathcal{I}
\quad\text{\textbf{open}}
\\
&\text{Stationary }
\mathbb{E}_\pi[\chi]\le -\chi_{\min}<0
\quad\text{\textbf{open (ACE)}}
\\
&\text{Bridge }
\Psi\to N_\star
\quad\text{\textbf{open; only after ACE}}
\\
&\text{No-Go stands until completed map + stationary mean.}
\end{aligned}
}
\end{equation}

% end
